\section{Problem Formulation} \label{sec:problem}
%\subsection{Problem Formulations}
We focus on classification problems in this paper. Let $\gZ = \gX \times \gY$, where $\gX$ denotes the input space, and $\gY=\{1,\cdots,K\}$ denotes the label space. Given a base-model training set $\mathcal{D}_B=\left\{\boldsymbol{x}^B_{i}, y^B_{i}\right\}_{i=1}^{N_B}\in \gZ^{N_B} $ containing i.i.d. samples generated from the distribution $P^{B}_{Z}$, a pretrained base model $\vh \circ \Phi:\gX \rightarrow \Delta^{K-1}$ is constructed, where $\Phi: \gX \rightarrow \sR^{l}$ and $\vh: \sR^{l} \rightarrow \Delta^{K-1}$ denote two complementary components of the neural network. More specifically,  $\Phi\left(\vx\right) = \Phi\left(\vx; \vw_{\phi}\right)$ stands for the intermediate feature representation of the base model, and the model output $\vh\left(\Phi(\vx)\right) = \vh\left(\Phi(\vx; \vw_{\phi}); \vw_h\right)$ denotes the  predicted label distribution $P_B\left( \vy \mid \Phi(\vx)\right) \in \Delta^{K-1}$ given input sample $\vx$, where $\left(\vw_{\phi}, \vw_h\right)\in \gW$ are the parameters of the pretrained base model.

The performance of the base model is evaluated by a non-negative loss function $\ell_B: \gW \times \gZ \to \sR_+$, e.g., cross entropy loss. Thus, a standard way to obtain the pretrained base model is by minimizing the empirical risk over $\mathcal{D}_B$, i.e., $\gL_B(\vh \circ \Phi,\gD_B) \triangleq \frac{1}{N_B}\sum_{i=1}^{N_B} \ell_B \left(\vh \circ \Phi, (\vx^B_i,y^B_i) \right)$.   

% Denote $\gW$ as the hypothesis class.

Although the well-trained deep base model is able to achieve good prediction accuracy, the output label distribution $P_B\left(\vy \mid \Phi(\vx)\right)$ is usually  unreliable for uncertainty quantification, i.e., it can be overconfident or poorly calibrated. Without retraining the model from scratch, we are interested in improving the uncertainty quantification performance in an efficient post-hoc manner. We utilize a meta-model $\vg : \sR^{l} \rightarrow \mathcal{\tilde{Y}}$ with parameter $\vw_{g}\in \gW_g$ building on top of the base model. The meta-model shares the same feature extractor from the base model and generates an output $\tilde{\vy} = \vg\left(\Phi(\vx); \vw_g\right)$, where $\tilde{\vy}\in \mathcal{\tilde{Y}}$ can take any form, e.g., a distribution over $\Delta^{K-1}$ or a scalar. Given a meta-model training set $\mathcal{D}_M=\left\{\boldsymbol{x}^M_{i}, y^M_{i}\right\}_{i=1}^{N_M}$ with i.i.d. samples from the distribution $P^{M}_{Z}$, our goal is to obtain the meta-model by optimizing a training objective $\gL_M(\vg \circ \Phi,\gD_M) \triangleq \frac{1}{N_M}\sum_{i=1}^{N_M} \ell_M \left(\vg \circ \Phi, (\vx^M_i,y^M_i) \right)$, where $\ell_M: \gW_g \times \gW_{\phi} \times \gZ \to \sR_+$ is the loss function for the meta-model. 

% $\vh \circ \Phi = \argmin_{\vh \circ \Phi} \gL_B(\vh \circ \Phi,\gD_B)$,  With shared base-model components $\Phi$ and dataset $\gD_M$ generated i.i.d from distribution $P^M_{Z}$

In the following, we formally introduce the post-hoc uncertainty learning problem using meta-model.
\begin{problem}\label{prob:general}[Post-hoc Uncertainty Learning by Meta-model] Given a base model $\vh \circ \Phi$ learned from the base-model training set $\gD_B$, the uncertainty learning problem by meta-model is to learn the function $\vg$ using the meta-model training set $\gD_M$ and the shared feature extractor $\Phi$, i.e.,
\begin{equation}
\vg^* = \argmin_{\vg} \gL_M(\vg \circ \Phi,\gD_M),
\end{equation}
such that the output from the meta-model $\tilde{\vy} = \vg\left(\Phi(x)\right)$ equipped with an uncertainty metric function $\vu : \mathcal{\tilde{Y}} \rightarrow \sR$ is able to generate a robust uncertainty score $\vu\left(\tilde{\vy}\right)$.
\end{problem}

%\subsection{Motivations} \label{sec:motivation}
%Problem~\ref{prob:general} offers a general and flexible framework for the post-hoc uncertainty learning problem. 
Next, the most critical questions are how the meta-model should use the information extracted from the pretrained base model, what kinds of uncertainty the meta-model should aim to quantify, and finally, how to train the meta-model appropriately. 

% In Section~\ref{sec:method}, we will specify the structure of the meta-model $\vg$, the loss function $\ell_M$ used for training the meta-model, and the uncertainty metric function $\vu$ in Problem~\ref{prob:general}.

% Next, the most important questions are what kind of information the meta-model should extract from the pretrained base-mode, what kinds of uncertainty should meta-model aim to quantify, and finally, how to appropriately train the meta-model. 



% In this paper, we provide a simple and effective to realize this framework. The design of our proposed meta-model uncertain learning method is based on three high-level insights: (1) different intermediate layers of the base model usually extract various levels of feature representation, e.g., from low-level to high-frequency. It's important to leverage the rich feature representation to achieve better uncertainty quantification performance. (2) In order to quantify different kinds of uncertainty, i.e., total uncertainty and epistemic uncertainty for various uncertainty quantification applications, we propose a Bayesian meta-model to parameterize a Dirichlet distribution as the conjugate prior distribution over label distribution. (3) The post-hoc setting explicitly focus on the uncertain learning instead of prediction, an validation strategy specifically designed for uncertainty learning that is similar to resolving over-fitting issue in supervised learning is crucial to achieve optimal uncertainty quantification performance.

